\documentclass[11pt,letterpaper]{article}

\usepackage[top=1in, bottom=1in, left=1in, right=1in, headheight=14pt]{geometry}
\usepackage{fontspec}
\setmainfont{DejaVu Serif}
\setsansfont{DejaVu Sans}
\setmonofont{DejaVu Sans Mono}

\usepackage{xcolor}
\usepackage{titlesec}
\usepackage{fancyhdr}
\usepackage{enumitem}
\usepackage{hyperref}
\usepackage{booktabs}
\usepackage{tabularx}
\usepackage{longtable}
\usepackage{colortbl}
\usepackage{multirow}
\usepackage{array}
\usepackage{parskip}
\usepackage{tcolorbox}
\tcbuselibrary{skins,breakable}
\usepackage{graphicx}
\usepackage{lastpage}

% ── Color Scheme: Prussian Blueprint ────────────────────────────
\definecolor{primary}{HTML}{0C2340}
\definecolor{secondary}{HTML}{7B3A10}
\definecolor{tertiary}{HTML}{005F87}
\definecolor{accent}{HTML}{B06000}
\definecolor{lightprimary}{HTML}{E4EBF2}
\definecolor{lightsecondary}{HTML}{F5EAE2}
\definecolor{lighttertiary}{HTML}{DDF0F7}
\definecolor{rowgray}{HTML}{F5F5F5}
\definecolor{bordergray}{HTML}{BDBDBD}
\definecolor{subtlegray}{HTML}{757575}
\definecolor{darktext}{HTML}{212121}
\definecolor{blueprintcyan}{HTML}{007BA7}
\definecolor{draftsienna}{HTML}{A0522D}
\definecolor{blockred}{HTML}{B71C1C}
\definecolor{gateyellow}{HTML}{E65100}
\definecolor{annogreen}{HTML}{1B5E20}

% ── Section Formatting ───────────────────────────────────────────
\titleformat{\section}
  {\Large\bfseries\sffamily\color{primary}}
  {\thesection}{1em}{}
  [\vspace{-0.5em}\textcolor{bordergray}{\rule{\textwidth}{0.4pt}}]
\titleformat{\subsection}
  {\large\bfseries\sffamily\color{secondary}}{\thesubsection}{1em}{}
\titleformat{\subsubsection}
  {\normalsize\bfseries\sffamily\color{tertiary}}{\thesubsubsection}{1em}{}
\titlespacing*{\section}{0pt}{1.5em}{0.8em}
\titlespacing*{\subsection}{0pt}{1.2em}{0.5em}
\titlespacing*{\subsubsection}{0pt}{0.8em}{0.3em}

% ── Custom Environments ──────────────────────────────────────────
\newcommand{\stageheader}[2]{%
  \noindent\colorbox{#2}{\makebox[\textwidth][l]{%
    \textcolor{white}{\bfseries\sffamily\hspace{6pt}#1\hspace{6pt}}}}%
  \vspace{0.5em}\par%
}
\newtcolorbox{asciibox}{
  colback=rowgray, colframe=bordergray,
  arc=1pt, boxrule=0.5pt,
  left=6pt, right=6pt, top=4pt, bottom=4pt,
  fontupper=\small\ttfamily
}
\newtcolorbox{quotebox}{
  colback=lightprimary, colframe=primary,
  arc=2pt, boxrule=0.5pt,
  left=12pt, right=12pt, top=8pt, bottom=8pt
}
\newtcolorbox{blueprintbox}{
  colback={blueprintcyan!8}, colframe=blueprintcyan,
  arc=2pt, boxrule=0.5pt,
  left=10pt, right=10pt, top=6pt, bottom=6pt
}
\newtcolorbox{safetybox}{
  colback={blockred!6}, colframe=blockred,
  arc=2pt, boxrule=0.7pt,
  left=10pt, right=10pt, top=6pt, bottom=6pt
}
\newcommand{\metafield}[2]{\textbf{\sffamily #1} #2\par}
\newcolumntype{L}[1]{>{\raggedright\arraybackslash}p{#1}}
\newcolumntype{C}[1]{>{\centering\arraybackslash}p{#1}}
\newcommand{\tableheader}[1]{\textbf{\sffamily\textcolor{white}{#1}}}
\newcommand{\blocktag}{\textcolor{blockred}{\textbf{\sffamily[BLOCK]}}}
\newcommand{\gatetag}{\textcolor{gateyellow}{\textbf{\sffamily[GATE]}}}
\newcommand{\annotag}{\textcolor{annogreen}{\textbf{\sffamily[ANNOTATE]}}}

% ── Headers and Footers ──────────────────────────────────────────
\pagestyle{fancy}
\fancyhf{}
\fancyhead[L]{\small\sffamily\textcolor{subtlegray}{So You Want to Be a Draftsman}}
\fancyhead[R]{\small\sffamily\textcolor{subtlegray}{GSD Mission Package}}
\fancyfoot[C]{\small\sffamily\textcolor{subtlegray}{Page \thepage\ of \pageref{LastPage}}}
\renewcommand{\headrulewidth}{0.4pt}
\renewcommand{\footrulewidth}{0pt}

\hypersetup{
  colorlinks=true, linkcolor=secondary, urlcolor=secondary,
  pdfauthor={GSD Ecosystem --- Tibsfox Inc.},
  pdftitle={So You Want to Be a Draftsman -- Mission Package},
  pdfsubject={Vision to Mission Pipeline},
}

% ─────────────────────────────────────────────────────────────────
\begin{document}

% ── TITLE PAGE ────────────────────────────────────────────────────
\thispagestyle{empty}
\begin{center}
\vspace*{2.5cm}
{\small\sffamily\textcolor{primary}{\textbf{GSD RESEARCH MISSION PACKAGE}}}
\vspace{0.3cm}

\textcolor{bordergray}{\rule{8cm}{0.4pt}}
\vspace{1.5cm}

{\fontsize{26}{32}\selectfont\bfseries\sffamily\textcolor{primary}{So You Want to Be\\[0.3em]a Draftsman}}

\vspace{0.8cm}
{\Large\sffamily\textcolor{secondary}{Blueprints, Technical Documentation,\\Simulation \& Testing Frameworks,\\Technical Drawing, Vector Graphics}}

\vspace{0.5cm}
{\large\sffamily\itshape\textcolor{tertiary}{Scientifically Accurate \& Mathematically Provable\\Technical Drawing for Real-World Construction}}

\vspace{2.5cm}
{\sffamily\textcolor{accent}{Vision $\rightarrow$ Research $\rightarrow$ Mission}}\\[0.3em]
{\small\sffamily\textcolor{accent}{Full Pipeline Execution}}

\vspace{2.5cm}
{\small\sffamily\textcolor{subtlegray}{Date: March 29, 2026 \quad|\quad Status: Ready for GSD Orchestrator}}\\[0.3em]
{\small\sffamily\textcolor{subtlegray}{Activation Profile: Fleet (17 roles) \quad|\quad Est.\ 6 parallel modules \quad|\quad $\sim$18 context windows}}
\end{center}
\newpage

% ── TABLE OF CONTENTS ─────────────────────────────────────────────
\tableofcontents
\newpage

% ═══════════════════════════════════════════════════════════════════
\stageheader{STAGE 1 --- VISION DOCUMENT}{primary}
% ═══════════════════════════════════════════════════════════════════

\section{So You Want to Be a Draftsman --- Vision Guide}

\metafield{Date:}{March 29, 2026}
\metafield{Status:}{Initial Vision / Research Complete}
\metafield{Depends on:}{gsd-skill-creator dev branch; NASA Mission Series Parts A--H templates}
\metafield{Context:}{Full deep-research mission covering technical drawing, blueprint creation, vector graphics, simulation frameworks, and prototype development pipelines within the GSD ecosystem.}

\subsection{Vision}

Before a bridge stands, before a spacecraft flies, before a maker's hands touch real metal --- there is a drawing. A line on vellum. A curve defined by a formula. A tolerance expressed as a symbol that says, precisely, how much error the universe is permitted to introduce before function becomes failure. The draftsman is the first builder. Everything else is execution of their geometry.

We live in an era when the distance between a drawing and a physical object has collapsed to nearly zero. A Python script generates an SVG. The SVG feeds a parametric model. The model drives a CNC router. Chips fall. A part emerges. The drawing and the object are the same thing, separated only by a coordinate transformation and some material constraints. This is not a metaphor --- it is the actual pipeline. The GSD ecosystem is built for this world: where the drawing is the program, the blueprint is the spec, and the draftsman's discipline is the engineering team's safety net.

``So You Want to Be a Draftsman'' is a meta deep-research mission that builds the full knowledge infrastructure for technically accurate, mathematically provable drawing and documentation within gsd-skill-creator. It does not merely explain how to draw --- it builds the skills, simulation pipelines, testing frameworks, and career pathways that make drawing a complete engineering discipline. The mission covers six interconnected domains: drawing standards and geometric dimensioning \& tolerancing (GD\&T); technical documentation and bill of materials systems; programmatic vector graphics and CNC-ready output; simulation and testing frameworks anchored in NASA-STD-7009B; materials science for real-world grounding; and the prototype development pipeline from drawing to physical object.

The Amiga Principle is alive in this domain. A draftsman who understands ASME Y14.5-2018 --- the complete symbolic language of manufacturing tolerances --- can specify, in a single feature control frame, requirements that would take paragraphs of prose to communicate ambiguously. That is architectural leverage: one symbol, unambiguous, machine-readable, inspection-provable, legally defensible, globally understood. The specialized execution path outperforms brute-force description every time.

\subsection{Problem Statement}

\begin{enumerate}
  \item \textbf{Drawings that cannot be built.} Technical drawings produced without knowledge of manufacturing constraints --- tolerances tighter than available processes, surface finishes unachievable with specified materials, assemblies with tolerance stackups that guarantee interference --- waste engineering hours and prototype budgets. A drawing that cannot be built is not a drawing; it is a hypothesis.
  \item \textbf{Standards fragmentation.} ASME Y14.5-2018 (US), ISO 1101:2017 (international), ISO 128-2:2020 (line types), ISO 2768 (general tolerances), NASA-STD-7009B (simulation credibility) --- these standards overlap, conflict, and update. Without a unified reference framework, teams default to informal conventions that break at every organizational boundary.
  \item \textbf{Simulation disconnected from drawing.} Finite element analysis and computational fluid dynamics are routinely performed on geometry that diverges from the drawing --- features simplified for meshing, tolerances ignored, material properties estimated. The simulation validates a model that does not match what will be manufactured. NASA-STD-7009B's credibility assessment framework exists precisely because this disconnection causes mission failures.
  \item \textbf{Vector graphics as an afterthought.} SVG is the most powerful 2D drawing format available: resolution-independent, mathematically precise, programmatically generated, web-native, and CNC-compatible. It is almost never taught alongside traditional drafting standards. The result is a generation of engineers who draw in expensive proprietary tools and cannot generate a compliant technical drawing from code.
  \item \textbf{No bridge from drawing to maker.} The pathway from a standards-compliant technical drawing to a physical prototype --- through DXF export, CAM toolpath generation, material selection, and fabrication method choice --- is treated as a separate discipline from drafting. For a maker, a student, or a small team, this gap means drawings stay on screen. The prototype never gets built.
\end{enumerate}

\subsection{Core Concept}

\textbf{Draw $\rightarrow$ Prove $\rightarrow$ Simulate $\rightarrow$ Build $\rightarrow$ Verify.}

Every stage of the pipeline has a mathematical guarantee. Drawings satisfy ASME Y14.5-2018 or ISO GPS. Simulations satisfy NASA-STD-7009B credibility criteria. Prototypes satisfy the original drawing's tolerance requirements as measured by inspection. The GSD draftsman does not guess --- they prove.

\subsubsection{Domain Architecture}

\begin{asciibox}
DRAWING STANDARDS LAYER
  ASME Y14.5-2018 (GD\&T) ---- ISO 128-2:2020 (Lines)
  ISO 1101:2017 (GPS)     ---- ISO 2768 (General Tol.)
  NASA-STD-7009B (M\&S)    ---- ANSI Y14.1 (Sheet Format)
              |
              v
PROGRAMMATIC DRAWING LAYER
  Python svgwrite/ezdxf -------- OpenSCAD (CSG Scripting)
  FreeCAD API (Python)  -------- LibreCAD (2D Drafting)
  CadQuery              -------- CodeToCAD (Multi-tool)
              |
              v
SIMULATION \& TESTING LAYER
  FEA (FreeCAD FEM / CalculiX) ---- CFD (OpenFOAM)
  NASA-STD-7009B Credibility   ---- ASME V\&V 40
  Tolerance Stackup Analysis   ---- Materials DB
              |
              v
PROTOTYPE PIPELINE LAYER
  DXF/SVG --> CAM (PyCAM) --> G-code --> CNC/Laser/3DP
  BOM generation --> Material procurement
  Inspection protocol --> CMM / manual gauge
              |
              v
GSD COLLEGE OF KNOWLEDGE DEPOSIT
  Skills: svg-technical-drawer, gdt-annotator,
          fea-pipeline, bom-generator, blueprint-generator,
          tolerance-calculator, maker-blueprint
  CK YAML entries in engineering department
\end{asciibox}

\subsection{Research Modules}

\subsubsection{Module 1: Drawing Standards and GD\&T}
\textbf{Purpose:} Build complete literacy in ASME Y14.5-2018 and ISO GPS --- symbols, rules, datums, tolerance zones, feature control frames, and their mathematical foundations.

\textbf{Key domains:} 14 geometric characteristic symbols (form, orientation, location, runout, profile); datum reference frames and degrees of freedom; material condition modifiers (MMC, LMC, RFS); tolerance stackup arithmetic; ASME Y14.5.1M mathematical definitions.

\textbf{Output:} GD\&T symbol reference guide (all 14 characteristics with worked examples); tolerance stackup calculator (algebraic and statistical methods); ASME vs.\ ISO comparison matrix; SKILL.md: \texttt{gdt-annotator}.

\subsubsection{Module 2: Technical Documentation and BOM Systems}
\textbf{Purpose:} Define the complete drawing package: title blocks, revision history, BOM structure, drawing tree relationships, notes conventions, and digital thread integration.

\textbf{Key domains:} ASME Y14.1 sheet formats (A--E sizes); title block fields (title, revision, material, finish, tolerances, projection angle); BOM hierarchy (top-level assembly $\to$ subassembly $\to$ part); DMS (Drawing Management System) schema; model-based definition (MBD) transition.

\textbf{Output:} Title block generator (Python + SVG); BOM schema (JSON); revision control pattern; SKILL.md: \texttt{bom-generator}.

\subsubsection{Module 3: Vector Graphics and Programmatic Drawing}
\textbf{Purpose:} Teach the full programmatic drawing pipeline: SVG generation from Python, DXF output for CAD interoperability, OpenSCAD parametric modeling, and CNC-ready G-code generation.

\textbf{Key domains:} SVG coordinate system and viewBox; \texttt{svgwrite} library; \texttt{ezdxf} for DXF generation; \texttt{CadQuery} for parametric parts; \texttt{PyCAM} and \texttt{dxf2gcode} for CAM; G-code structure (rapid/linear move, tool change, coordinate systems).

\textbf{Output:} Programmatic drawing cookbook (20 drawing patterns); CNC pipeline spec; SVG-to-DXF converter; SKILL.md: \texttt{svg-technical-drawer}, \texttt{blueprint-generator}.

\subsubsection{Module 4: Simulation and Testing Frameworks}
\textbf{Purpose:} Build the simulation pipeline from drawing geometry through FEA/CFD to NASA-STD-7009B credibility assessment, with open-source tools (FreeCAD FEM, CalculiX, OpenFOAM).

\textbf{Key domains:} Finite element method fundamentals (discretization, mesh convergence, boundary conditions, material properties); FreeCAD FEM workbench (CalculiX solver interface); OpenFOAM CFD (mesh generation, boundary conditions, turbulence models); NASA-STD-7009B credibility levels 1--4; ASME V\&V 40 verification and validation; uncertainty quantification.

\textbf{Output:} FEA pipeline tutorial (bracket under load, thermal gradient); CFD pipeline tutorial (duct flow); NASA-STD-7009B credibility worksheet template; SKILL.md: \texttt{fea-pipeline}.

\subsubsection{Module 5: Materials Science Reference}
\textbf{Purpose:} Provide the physical constants, material properties, and fabrication constraints that make drawings buildable --- connecting tolerances to process capabilities and material behavior.

\textbf{Key domains:} Common engineering materials (aluminum alloys, steel grades, HDPE, carbon fiber); material property tables (tensile strength, yield strength, Young's modulus, thermal expansion coefficient, density); manufacturing process capability (Cpk, IT grades per ISO 286-1); surface finish Ra values by process; fastener standards (ISO 898, SAE J429).

\textbf{Output:} Materials reference table (12 materials $\times$ 10 properties); process capability map (IT grades vs.\ manufacturing process); SKILL.md: \texttt{tolerance-calculator}.

\subsubsection{Module 6: Prototype Development Pipeline}
\textbf{Purpose:} Close the loop from technical drawing to physical object: DXF export, maker tool selection (laser cutter, 3D printer, CNC router, hand tools), BOM-to-procurement, and inspection protocol.

\textbf{Key domains:} Drawing export formats (DXF, SVG, STEP, STL, G-code); maker tool capability comparison; low-cost fabrication methods (foam board, plywood, aluminum sheet, 3D print); inspection tools (calipers, gauge blocks, go/no-go gauges, laser measurement); DIY builds: L1 \$15 (hand tools), L2 \$80 (3D printer), L3 \$200 (laser cutter), L4 \$600 (CNC router).

\textbf{Output:} Maker tool selection matrix; inspection protocol template; prototype BOM sheet; SKILL.md: \texttt{maker-blueprint}.

\subsection{Chipset Configuration}

\begin{asciibox}
name: so-you-want-to-be-a-draftsman
version: 1.0.0

agnus:    \# Opus ~30%
  role: FLIGHT + SYNTHESIS
  responsibility: >
    GD\&T mathematical correctness review;
    NASA-STD-7009B credibility assessment;
    module integration; career pathway synthesis;
    safety warden for tolerance errors

denise:   \# Sonnet ~60%
  role: EXEC-A through EXEC-F
  responsibility: >
    All six module implementations;
    drawing standards documentation;
    programmatic drawing cookbook;
    FEA/CFD pipeline tutorials;
    materials reference tables;
    prototype pipeline

paula:    \# Sonnet ~60%
  role: PUBLICATION + CK
  responsibility: >
    SKILL.md file generation (7 skills);
    CK YAML deposit to engineering department;
    PDF compilation; index.html generation

gary:     \# Haiku ~10%
  role: SCAFFOLD + SCHEMA
  responsibility: >
    Directory structure; config files;
    JSON schemas; BOM template scaffolds;
    title block field definitions

model\_allocation:
  opus: 30%
  sonnet: 60%
  haiku: 10%
\end{asciibox}

\subsection{Scope Boundaries}

\textbf{In Scope (v1.0):}
\begin{itemize}[noitemsep]
  \item ASME Y14.5-2018 (US standard) and ISO 1101:2017 / ISO 128-2:2020 (international)
  \item 2D technical drawing (orthographic projection, sections, details)
  \item Programmatic drawing: Python \texttt{svgwrite}, \texttt{ezdxf}, \texttt{CadQuery}
  \item Open-source simulation: FreeCAD FEM, CalculiX, OpenFOAM
  \item NASA-STD-7009B credibility levels 1--4 for simulation validation
  \item Prototype methods: 3D printing, laser cutting, CNC routing, hand fabrication
  \item Seven SKILL.md files deposited to College of Knowledge engineering department
\end{itemize}

\textbf{Out of Scope:}
\begin{itemize}[noitemsep]
  \item Commercial CAD (SolidWorks, CATIA, NX, Revit) --- proprietary, not ecosystem-aligned
  \item Structural engineering codes (AISC, ACI, Eurocode) --- separate mission required
  \item Electrical / PCB schematics --- covered in future Electronics mission
  \item Architectural drawings (IFC, BIM workflows) --- separate AEC mission
  \item 3D parametric modeling (FreeCAD Part Design) as primary focus --- treated as reference tool
\end{itemize}

\subsection{Success Criteria}

\begin{enumerate}
  \item All 14 GD\&T geometric characteristic symbols are documented with definition, symbol, zone description, and worked example per ASME Y14.5-2018.
  \item Tolerance stackup calculator produces correct results (algebraic and RSS statistical) verified against hand-calculated reference cases.
  \item Programmatic SVG drawing pipeline produces a standards-compliant title block, orthographic three-view drawing, and dimension annotations from Python code alone.
  \item DXF export from Python (\texttt{ezdxf}) is importable by FreeCAD and LibreCAD without errors; dimensions, layers, and line types preserved.
  \item FEA tutorial (cantilever bracket, 1018 steel, point load) converges within 2\% of analytical beam theory solution using FreeCAD FEM / CalculiX.
  \item NASA-STD-7009B credibility worksheet is complete, field-by-field, with definitions and a worked example at Credibility Level 2.
  \item Materials reference table covers 12 materials with: density, Young's modulus, yield strength, UTS, thermal expansion coefficient, and typical manufacturing process.
  \item Prototype pipeline tutorial guides a student from a completed SVG drawing to G-code (PyCAM) for CNC routing without commercial software.
  \item All seven SKILL.md files (\texttt{svg-technical-drawer}, \texttt{blueprint-generator}, \texttt{gdt-annotator}, \texttt{bom-generator}, \texttt{tolerance-calculator}, \texttt{fea-pipeline}, \texttt{maker-blueprint}) are deployed to \texttt{/mnt/skills/user/} and validate against the SKILL.md schema.
  \item GSD Try Sessions exist for each module: hands-on activities completable with household materials (graph paper, ruler, compass, cardboard) or free software.
  \item ASME vs.\ ISO comparison matrix covers the 10 highest-impact difference areas with concrete example drawings.
  \item Career pathway section documents four tracked pathways: Mechanical Draftsman, Simulation Engineer, Maker/Fabricator, Technical Writer --- with certifications, salary ranges (BLS data), and GSD skill-to-career mapping.
\end{enumerate}

\subsection{Relationship to Other Documents}

\begin{center}
\rowcolors{2}{rowgray}{white}
\begin{tabularx}{\textwidth}{L{4cm}L{4cm}X}
\toprule
\rowcolor{primary}
\tableheader{Document} & \tableheader{Relationship} & \tableheader{Data Flow} \\
\midrule
NASA Part B (Single Mission) & Consumer & Part B science data feeds materials reference and FEA case studies \\
NASA Part D (Engineering) & Sibling & Part D engineering math workbook cross-references GD\&T and tolerance content \\
NASA Part E (Science+Skills) & Consumer & Part E skill expansion uses \texttt{svg-technical-drawer} and \texttt{fea-pipeline} \\
NASA Part G (Refinement) & Consumer & Part G TSPB chapters pull from materials reference and drawing standards tables \\
NASA Part H Template & Sibling & Part H template is the per-mission application; this mission is the knowledge infrastructure \\
The Space Between (TSPB) & Feeder & Materials properties, tolerance math, FEA derivations feed TSPB Layer 4 (Vector Calculus) \\
\bottomrule
\end{tabularx}
\end{center}

\subsection{The Through-Line}

The draftsman is the person who lives in the spaces between design intent and physical reality. The engineer imagines a force-bearing structure; the draftsman specifies the datum reference frame that tells the machinist where on the part to begin measuring. The scientist models a fluid; the draftsman draws the inlet geometry that determines whether the model is even solving the right equations. The maker has an idea; the draftsman produces the blueprint that transforms the idea into a set of instructions that survive the gap between one person's hands and another's.

This is the GSD principle in its most literal form. The ecosystem exists to reduce friction between creative intent and realized outcome. The draftsman's discipline --- standards literacy, mathematical precision, simulation validation, prototype-to-inspection closure --- is that friction reduction made visible. Every symbol in a feature control frame is a conversation compressed to a glyph. Every mesh convergence test is a question answered before the prototype is built. Every materials table is a physical law made navigable.

``So You Want to Be a Draftsman'' teaches not just drawing but the complete epistemic pipeline from idea to provable object.

% ═══════════════════════════════════════════════════════════════════
\newpage
\stageheader{STAGE 2 --- RESEARCH REFERENCE}{secondary}
% ═══════════════════════════════════════════════════════════════════

\section{Research Reference}

\subsection{How to Use This Document}

This reference compiles authoritative findings across all six modules. Each section is organized for rapid lookup during mission execution. All numerical claims are sourced; unsourced approximations are marked \textit{(est.)}.

\textbf{Key Source Organizations:}
\begin{itemize}[noitemsep]
  \item ASME (American Society of Mechanical Engineers) --- Y14.5-2018, Y14.1, Y14.5.1M
  \item ISO (International Organization for Standardization) --- 128-2:2020, 1101:2017, 2768, 286-1
  \item NASA Office of the Chief Engineer --- NASA-STD-7009B (2024), NASA-HDBK-7009A
  \item NIST (National Institute of Standards and Technology) --- material property data, measurement standards
  \item FreeCAD Project / FEM Workbench --- open-source FEA documentation
  \item Ansys, Siemens, Autodesk --- FEA educational documentation (public materials only)
  \item BLS (Bureau of Labor Statistics) --- career salary and employment data
\end{itemize}

\subsection{Module 1: Drawing Standards and GD\&T}

\subsubsection{ASME Y14.5-2018 Overview}

<ASME Y14.5-2018 (R2024)> establishes symbols, rules, definitions, requirements, defaults, and recommended practices for stating and interpreting geometric dimensioning and tolerancing. It is the primary GD\&T standard in the United States, Canada, and Australia. <ISO 1101:2017> is the international equivalent; the two standards share approximately 65--70\% agreement on symbols but diverge significantly on fundamental interpretation principles (axis method vs.\ surface boundary method).

\begin{center}
\rowcolors{2}{rowgray}{white}
\begin{tabularx}{\textwidth}{L{2.8cm}C{1.4cm}L{2.5cm}X}
\toprule
\rowcolor{primary}
\tableheader{Characteristic} & \tableheader{Symbol} & \tableheader{Zone Type} & \tableheader{Datum Required?} \\
\midrule
Flatness & \texttt{[]=} & Two parallel planes & No \\
Circularity (Roundness) & \texttt{O} & Two concentric circles & No \\
Cylindricity & \textit{cyl} & Coaxial cylinders & No \\
Straightness & --- & Two parallel lines & No \\
Profile of a Line & \texttt{\textasciicircum} & Two offset lines & Optional \\
Profile of a Surface & \texttt{\^{}} & Two offset surfaces & Optional \\
Angularity & $\angle$ & Two parallel planes & Yes \\
Perpendicularity & $\perp$ & Two parallel planes & Yes \\
Parallelism & $\parallel$ & Two parallel planes & Yes \\
Position & $\oplus$ & Cylinder/sphere & Yes \\
Concentricity (removed 2018) & --- & --- & --- \\
Symmetry (removed 2018) & --- & --- & --- \\
Circular Runout & \texttt{/--} & Circle in rotation & Yes \\
Total Runout & $\Rightarrow$ & Cylinder in rotation & Yes \\
\bottomrule
\end{tabularx}
\end{center}

\textbf{Key 2018 Changes} (from ASME Y14.5-2009): Concentricity and symmetry symbols eliminated; profile of a surface now preferred for location tolerances; plus/minus location tolerances moved to Appendix (likely removed in next revision); model-based definition (MBD) integration added throughout; composite position tolerance clarified.

\subsubsection{Tolerance Stackup: Algebraic and RSS}

For an assembly with $n$ independent tolerances $t_1, t_2, \ldots, t_n$:

\textbf{Worst-case (algebraic):} $T_{assy} = \sum_{i=1}^{n} |c_i \cdot t_i|$ where $c_i$ is the sensitivity factor (typically $\pm 1$).

\textbf{Statistical (RSS):} $T_{assy} = \sqrt{\sum_{i=1}^{n} (c_i \cdot t_i)^2}$ --- valid when tolerances are independent, processes are centered, and $n \geq 4$.

Typical IT grade tolerances (ISO 286-1) for a 25 mm nominal dimension:
\begin{center}
\rowcolors{2}{rowgray}{white}
\begin{tabularx}{\textwidth}{C{2cm}C{2cm}C{2.5cm}X}
\toprule
\rowcolor{primary}
\tableheader{IT Grade} & \tableheader{Tol. ($\mu$m)} & \tableheader{Tol. (in)} & \tableheader{Typical Process} \\
\midrule
IT4 & 9 & 0.00035 & Precision grinding \\
IT6 & 13 & 0.00051 & Finish turning, broaching \\
IT7 & 21 & 0.00083 & Turning, boring, reaming \\
IT9 & 52 & 0.00205 & Cold drawing, extrusion \\
IT11 & 130 & 0.00512 & Rough turning, drilling \\
IT14 & 520 & 0.02047 & Sand casting, hot rolling \\
\bottomrule
\end{tabularx}
\end{center}

\subsubsection{ASME vs.\ ISO: 10 High-Impact Differences}

\begin{center}
\rowcolors{2}{rowgray}{white}
\begin{tabularx}{\textwidth}{L{4cm}L{4.5cm}X}
\toprule
\rowcolor{primary}
\tableheader{Topic} & \tableheader{ASME Y14.5-2018} & \tableheader{ISO GPS} \\
\midrule
Default interpretation & Axis/centerplane method & Surface boundary method \\
Size rule (Rule \#1) & Envelope principle mandatory & Must be explicitly invoked \\
Concentricity & Removed (2018) & Still present (ISO 1101) \\
Symmetry & Removed (2018) & Still present \\
Position MMC & Bonus tolerance for mating & Virtual condition boundary \\
Surface texture & Separate (ASME B46.1) & ISO 1302 / ISO 4287 \\
Projection angle & Third-angle default & First-angle default \\
Drawing sheet sizes & ANSI (A--E, 8.5x11 to 34x44) & ISO 216 (A0--A4) \\
Profile without datum & Controls form only & Controls form and orientation \\
Datum feature symbol & Open square on extension line & Filled triangle on line \\
\bottomrule
\end{tabularx}
\end{center}

\subsection{Module 2: Technical Documentation and BOM Systems}

\subsubsection{ISO 128 Line Type Standard (2020 Revision)}

ISO 128-2:2020 consolidated six former parts into a unified line standard. It specifies line types, designations, configurations, and general drafting rules. Two line widths are specified for mechanical drawings: wide lines ($\approx 0.5$--$0.7$ mm) and narrow lines ($\approx 0.25$--$0.35$ mm) in a 2:1 ratio. ISO 128-3 governs views, sections, and cuts (including first-angle and third-angle projection conventions).

Key line types for technical drawings:
\begin{center}
\rowcolors{2}{rowgray}{white}
\begin{tabularx}{\textwidth}{L{3cm}L{3cm}L{2.5cm}X}
\toprule
\rowcolor{primary}
\tableheader{Line Name} & \tableheader{ISO 128 Type} & \tableheader{Width} & \tableheader{Application} \\
\midrule
Visible outline & 01.2 continuous wide & Wide & Visible edges of object \\
Hidden outline & 02.1 dashed narrow & Narrow & Edges not visible in view \\
Center line & 04.1 long-dash dotted narrow & Narrow & Axes of symmetry \\
Cutting plane & 04.2 long-dash double-dot & Wide ends & Section cut indication \\
Break line (long) & 01.1 continuous narrow freehand & Narrow & Long break indication \\
Dimension line & 01.1 continuous narrow & Narrow & Dimension terminations \\
Extension line & 01.1 continuous narrow & Narrow & Dimension extension \\
\bottomrule
\end{tabularx}
\end{center}

\subsubsection{Drawing Package BOM Schema}

A complete drawing package includes: title drawing (assembly level), subassembly drawings, detail drawings (individual parts), and the BOM. The BOM links drawing numbers to parts via balloon callouts.

\begin{asciibox}
BOM JSON Schema (v1.0):
{
  "bom\_id": "BOM-2026-001",
  "assembly\_drawing": "DWG-2026-001-A",
  "revision": "B",
  "items": [
    {
      "item\_no": 1,
      "part\_number": "PN-2026-001-001",
      "description": "Base Plate",
      "qty": 1,
      "material": "6061-T6 Aluminum",
      "finish": "Anodize Type II",
      "source": "MAKE",
      "mass\_kg": 0.145,
      "drawing": "DWG-2026-001-001"
    },
    ...
  ],
  "total\_mass\_kg": 2.340,
  "notes": "All dimensions in mm unless noted"
}
\end{asciibox}

\subsection{Module 3: Vector Graphics and Programmatic Drawing}

\subsubsection{SVG Coordinate System for Technical Drawing}

SVG uses a left-handed coordinate system with the origin at the top-left and Y increasing downward. For technical drawings, it is conventional to establish a working coordinate system with $+Y$ pointing up via a viewport transform: \texttt{transform="translate(0, viewbox\_height) scale(1,-1)"}.

Key SVG drawing elements for technical drawing:
\begin{itemize}[noitemsep]
  \item \texttt{<line>} --- straight lines (visible, hidden, center, dimension)
  \item \texttt{<circle>}, \texttt{<ellipse>} --- circular features
  \item \texttt{<path>} --- complex curves, arcs, splines
  \item \texttt{<text>} --- dimension values, notes, title block fields
  \item \texttt{<marker>} --- arrowheads for dimension terminations
  \item \texttt{<symbol>} / \texttt{<use>} --- GD\&T symbols as reusable components
\end{itemize}

\subsubsection{Open-Source Drawing Tool Ecosystem (2026)}

\begin{center}
\rowcolors{2}{rowgray}{white}
\begin{tabularx}{\textwidth}{L{2.5cm}L{2cm}L{2.5cm}X}
\toprule
\rowcolor{primary}
\tableheader{Tool} & \tableheader{Type} & \tableheader{License} & \tableheader{Best For} \\
\midrule
FreeCAD (1.0+) & 3D parametric + FEA & LGPL-2.0 & Parametric design, FEM, CAM workbenches, Python scripting \\
LibreCAD & 2D drafting & GPL-2.0 & Standards-compliant 2D drawings, DXF output \\
OpenSCAD & Code-based CSG & GPL-2.0 & Parametric parts from scripts; 3D print ready \\
CadQuery & Python CSG & Apache-2.0 & Python-first parametric 3D, STEP/DXF/SVG export \\
svgwrite & Python SVG & MIT & Programmatic SVG generation \\
ezdxf & Python DXF & MIT & DXF read/write, dimension entities, layers \\
PyCAM & Python CAM & GPL-3.0 & STL/DXF to G-code; CNC toolpath generation \\
dxf2gcode & Python CAM & GPL-3.0 & 2D DXF to G-code; grbl-compatible \\
CalculiX & FEA solver & GPL-2.0 & FreeCAD FEM backend; structural/thermal analysis \\
OpenFOAM & CFD solver & GPL-3.0 & Fluid dynamics simulation \\
\bottomrule
\end{tabularx}
\end{center}

\subsection{Module 4: Simulation and Testing Frameworks}

\subsubsection{NASA-STD-7009B Credibility Assessment}

NASA-STD-7009B (2024 revision, replacing 7009A) establishes standard requirements for Models and Simulations (M\&S) used in NASA engineering decisions. The Credibility Assessment Scale (CAS) rates M\&S from Level 0 (no evidence) to Level 4 (validated against flight data). Key components of credibility:

\begin{center}
\rowcolors{2}{rowgray}{white}
\begin{tabularx}{\textwidth}{L{3cm}X}
\toprule
\rowcolor{primary}
\tableheader{CAS Factor} & \tableheader{Definition} \\
\midrule
Verification & Does the code solve the equations correctly? (code verification, solution verification) \\
Validation & Do the simulation results match physical reality? (comparison to experiment) \\
Input Pedigree & Are boundary conditions, material properties, and geometry representative of the real system? \\
Results Uncertainty & Is uncertainty in simulation outputs quantified (epistemic and aleatory)? \\
Use Exposure & Has the M\&S been used in similar applications before? \\
Technical Review & Has the M\&S been peer-reviewed? \\
\bottomrule
\end{tabularx}
\end{center}

Level 4 (highest) requires experimental validation data with the same geometry, loads, and environment as the operational system. Level 2 (achievable for prototype work) requires comparison to referent data from expert opinion or other M\&S plus rigorous uncertainty characterization.

\subsubsection{FEA Pipeline: Cantilever Bracket Reference Case}

For a cantilever beam of length $L = 0.2$ m, width $b = 0.02$ m, height $h = 0.01$ m, load $P = 100$ N at tip, material 1018 steel ($E = 200$ GPa, $\nu = 0.29$, $\rho = 7870$ kg/m\textsuperscript{3}):

\textbf{Analytical tip deflection:} $\delta = \frac{PL^3}{3EI}$ where $I = \frac{bh^3}{12} = 1.667 \times 10^{-9}$ m\textsuperscript{4}

$\delta = \frac{100 \times (0.2)^3}{3 \times 200 \times 10^9 \times 1.667 \times 10^{-9}} = \frac{0.8}{1000.2} \approx 0.800$ mm

\textbf{Maximum bending stress:} $\sigma_{max} = \frac{PL \cdot c}{I} = \frac{100 \times 0.2 \times 0.005}{1.667 \times 10^{-9}} = 60.0$ MPa $\ll \sigma_y = 370$ MPa for 1018 steel.

FEA convergence criterion: mesh-refined solution within 2\% of $\delta_{analytical}$ and $\sigma_{analytical}$ before proceeding to design iterations.

\subsection{Module 5: Materials Science Reference}

\begin{center}
\rowcolors{2}{rowgray}{white}
\begin{tabularx}{\textwidth}{L{2.8cm}C{1.2cm}C{1.3cm}C{1.3cm}C{1.3cm}C{1.3cm}}
\toprule
\rowcolor{primary}
\tableheader{Material} & \tableheader{Density (kg/m\textsuperscript{3})} & \tableheader{E (GPa)} & \tableheader{Yield (MPa)} & \tableheader{UTS (MPa)} & \tableheader{CTE ($\mu$m/m\textperiodcentered K)} \\
\midrule
1018 Steel & 7870 & 200 & 370 & 440 & 11.7 \\
4140 Steel (HT) & 7850 & 205 & 655 & 1020 & 12.3 \\
304 Stainless & 8000 & 193 & 215 & 505 & 17.2 \\
6061-T6 Alum. & 2700 & 69 & 276 & 310 & 23.6 \\
7075-T6 Alum. & 2810 & 72 & 503 & 572 & 23.4 \\
Grade 2 Ti & 4510 & 105 & 275 & 345 & 8.6 \\
HDPE & 955 & 1.1 & 26 & 37 & 120 \\
PLA (3D print) & 1240 & 3.5 & 50 & 65 & 68 \\
PETG (3D print) & 1270 & 2.1 & 43 & 53 & 80 \\
Carbon Fiber (UD) & 1600 & 70--130 & --- & 1500 & 0.0--1.0 \\
Plywood (birch) & 680 & 11 & --- & 40 & 4--10 \\
Borosilicate Glass & 2230 & 63 & --- & 49 & 3.3 \\
\bottomrule
\end{tabularx}
\end{center}

\textit{Sources: NIST Material Measurement Laboratory; MatWeb; ASM International Material Property Database.}

\subsection{Module 6: Prototype Development Pipeline}

\subsubsection{Maker Tool Capability Matrix}

\begin{center}
\rowcolors{2}{rowgray}{white}
\begin{tabularx}{\textwidth}{L{2.5cm}C{1.5cm}C{1.5cm}L{2.5cm}X}
\toprule
\rowcolor{primary}
\tableheader{Tool} & \tableheader{Cost (USD)} & \tableheader{Accuracy} & \tableheader{Materials} & \tableheader{CAD/CAM Input} \\
\midrule
Hand tools only & \$15 & $\pm$1 mm & Wood, cardboard, foam & Template / printout \\
FDM 3D printer & \$200--800 & $\pm$0.2 mm & PLA, PETG, ABS, TPU & STL / G-code \\
Laser cutter (K40) & \$400--600 & $\pm$0.1 mm & Acrylic, plywood, cardboard & SVG / DXF \\
Desktop CNC router & \$600--1500 & $\pm$0.05 mm & Wood, aluminum, PCB & G-code (PyCAM) \\
Drill press + mill & \$800--2500 & $\pm$0.02 mm & Metal, plastic & Manual from drawing \\
Pro CNC (Tormach) & \$8000+ & $\pm$0.005 mm & All metals & CAM software \\
\bottomrule
\end{tabularx}
\end{center}

\subsubsection{Drawing-to-Prototype Pipeline}

\begin{asciibox}
1. DRAW: Python (svgwrite / ezdxf) --> SVG / DXF
   - Apply GD\&T annotations
   - Specify material and surface finish in title block
   - Complete BOM

2. VERIFY: Tolerance stackup check
   - All mating features: verify MMC conditions
   - All dimensions: trace back to datum reference frame
   - FEA if structural (CalculiX via FreeCAD FEM)

3. EXPORT:
   - 3D print --> STL (FreeCAD / OpenSCAD export)
   - Laser cut --> SVG/DXF (direct use)
   - CNC router --> DXF --> PyCAM --> G-code
   - Sheet metal --> DXF --> laser/waterjet quote

4. FABRICATE:
   - Set material datum per drawing (align with DRF)
   - Machine/print/cut to drawing

5. INSPECT:
   - Calipers: linear dimensions (+-0.02 mm resolution)
   - Go/No-Go gauge: controlled features
   - Surface plate + height gauge: flatness, squareness
   - CMM if available: full GD\&T verification

6. DOCUMENT:
   - Update BOM with actual lot/batch numbers
   - Record measured values in inspection report
   - Note any NCRs (non-conformances)
\end{asciibox}

\subsection{Safety and Sensitivity Considerations}

\begin{center}
\rowcolors{2}{rowgray}{white}
\begin{tabularx}{\textwidth}{L{3cm}L{2cm}X}
\toprule
\rowcolor{primary}
\tableheader{Rule} & \tableheader{Type} & \tableheader{Application} \\
\midrule
Tolerance math must be verified analytically & \blocktag & All tolerance stackup calculations must have a hand-computed check case \\
FEA credibility per NASA-STD-7009B must be stated & \gatetag & All simulation results must declare their credibility level before use in design decisions \\
Material properties must cite source & \gatetag & All material property values used in FEA must reference NIST, MatWeb, or ASM International \\
DIY builds must include safety notes & \annotag & Laser cutter, CNC router, and power tool activities must include PPE requirements \\
Never specify tolerances tighter than process can achieve & \blocktag & A drawing with IT4 tolerances specified for sand casting is invalid; reject before fabrication \\
\bottomrule
\end{tabularx}
\end{center}

\subsection{Source Bibliography}

\textbf{Standards:}
\begin{itemize}[noitemsep]
  \item ASME Y14.5-2018 (R2024). \textit{Dimensioning and Tolerancing}. ASME, New York.
  \item ASME Y14.5.1M-1994. \textit{Mathematical Definition of Dimensioning and Tolerancing Principles}. ASME.
  \item ISO 128-1:2020. \textit{Technical Product Documentation --- General Principles of Representation}. ISO.
  \item ISO 128-2:2020. \textit{Basic Conventions for Lines}. ISO.
  \item ISO 128-3:2022. \textit{Views, Sections, and Cuts}. ISO.
  \item ISO 1101:2017. \textit{Geometrical Product Specifications --- Geometrical Tolerancing}. ISO.
  \item ISO 2768-1:1989. \textit{General Tolerances --- Linear and Angular Dimensions}. ISO.
  \item ISO 286-1:2010. \textit{ISO Code System for Tolerances on Linear Sizes}. ISO.
  \item NASA-STD-7009B (2024). \textit{Standard for Models and Simulations}. NASA Office of the Chief Engineer.
  \item NASA-HDBK-7009A (2016). \textit{Handbook for Models and Simulations}. NASA.
\end{itemize}

\textbf{Open-Source Software:}
\begin{itemize}[noitemsep]
  \item FreeCAD Project. \textit{FreeCAD 1.0 Documentation} (LGPL-2.0). \url{https://www.freecad.org}
  \item LibreCAD. \textit{LibreCAD 2.x Documentation} (GPL-2.0). \url{https://librecad.org}
  \item OpenFOAM Foundation. \textit{OpenFOAM v12 User Guide} (GPL-3.0). \url{https://openfoam.org}
  \item ezdxf library (MIT). \url{https://github.com/mozman/ezdxf}
  \item svgwrite library (MIT). \url{https://github.com/mozman/svgwrite}
  \item CadQuery (Apache-2.0). \url{https://github.com/CadQuery/cadquery}
\end{itemize}

% ═══════════════════════════════════════════════════════════════════
\newpage
\stageheader{STAGE 3 --- MISSION PACKAGE}{tertiary}
% ═══════════════════════════════════════════════════════════════════

\section{Milestone Specification}

\textbf{Mission Objective:} Build the complete GSD draftsman knowledge infrastructure: six research modules covering drawing standards, technical documentation, programmatic vector graphics, simulation frameworks, materials science, and the prototype pipeline --- plus seven SKILL.md deployments to the College of Knowledge engineering department --- within a single Fleet-profile execution.

\begin{center}
\rowcolors{2}{rowgray}{white}
\begin{tabularx}{\textwidth}{C{0.5cm}L{3.5cm}L{3cm}X}
\toprule
\rowcolor{primary}
\tableheader{\#} & \tableheader{Deliverable} & \tableheader{Format} & \tableheader{Acceptance Criterion} \\
\midrule
1 & GD\&T Symbol Reference & Markdown + SVG & All 14 characteristics, worked examples, ASME vs.\ ISO comparison \\
2 & Tolerance Stackup Calculator & Python script & Algebraic and RSS; verified against reference cases \\
3 & Programmatic Drawing Cookbook & Markdown + Python & 20 drawing patterns; compliant title block from code alone \\
4 & SVG/DXF Pipeline Tutorial & Markdown + Python & DXF importable by FreeCAD; dimensions/layers/linetypes preserved \\
5 & FEA Tutorial (Cantilever) & Markdown + FreeCAD & Within 2\% of analytical solution; CalculiX solver \\
6 & NASA-STD-7009B Worksheet & Markdown template & All CAS fields documented; Level 2 worked example \\
7 & Materials Reference Table & Markdown + JSON & 12 materials $\times$ 10 properties; sourced values \\
8 & Prototype Pipeline Tutorial & Markdown + code & SVG to G-code (PyCAM) without commercial software \\
9 & 7 $\times$ SKILL.md files & SKILL.md schema & Deployed to \texttt{/mnt/skills/user/}; schema-valid \\
10 & 4 Career Pathway Guides & Markdown & BLS salary ranges; certifications; GSD skill mapping \\
11 & GSD Try Sessions (6) & Markdown & One per module; household materials or free software \\
12 & ASME vs.\ ISO Matrix & Markdown table & 10 high-impact differences with example drawings \\
\bottomrule
\end{tabularx}
\end{center}

\section{Wave Execution Plan}

\subsection{Wave Summary}

\begin{center}
\rowcolors{2}{rowgray}{white}
\begin{tabularx}{\textwidth}{C{1.2cm}L{3cm}C{2cm}C{1.8cm}C{1.8cm}X}
\toprule
\rowcolor{primary}
\tableheader{Wave} & \tableheader{Name} & \tableheader{Mode} & \tableheader{Model} & \tableheader{Tokens} & \tableheader{Output} \\
\midrule
W0 & Foundation & Serial & Haiku & 15K & Shared schemas, BOM JSON, SKILL.md schema, title block template \\
W1 & Research Modules & Parallel (6) & Sonnet & 180K & 6 module research docs + draft SKILL.md files \\
W2 & Synthesis & Serial & Opus & 45K & Career pathways, ASME/ISO matrix, Try Sessions, integration \\
W3 & Publication & Serial & Sonnet & 40K & SKILL.md deployment, CK YAML, prototype pipeline, index.html \\
\bottomrule
\end{tabularx}
\end{center}

\textbf{Total estimated tokens:} $\approx$280K \quad \textbf{Context windows:} $\approx$18 \quad \textbf{Activation profile:} Fleet (17 roles)

\subsection{Wave 0: Foundation (Serial, Haiku, $<$5 min)}

\begin{center}
\rowcolors{2}{rowgray}{white}
\begin{tabularx}{\textwidth}{L{3.5cm}L{3cm}X}
\toprule
\rowcolor{primary}
\tableheader{Task} & \tableheader{Agent} & \tableheader{Output} \\
\midrule
Shared schemas & SCAFFOLD (Haiku) & BOM JSON schema v1.0; SKILL.md required fields; IT grade lookup table \\
Title block template & SCAFFOLD (Haiku) & SVG title block scaffold (ANSI A--D formats); field definitions \\
Drawing standards index & SCAFFOLD (Haiku) & Standard cross-reference: Y14.5-2018, ISO 128-2:2020, ISO 1101:2017, Y14.1 \\
GD\&T symbol index & SCAFFOLD (Haiku) & 14-characteristic index; zone type for each; datum required flag \\
\bottomrule
\end{tabularx}
\end{center}

\subsection{Wave 1: Parallel Research Tracks (6 concurrent, Sonnet)}

\textbf{Track A --- Drawing Standards and GD\&T} (EXEC-A, Sonnet, $\sim$32K tokens)\\
Input: GD\&T symbol index (W0). Output: complete GD\&T reference guide; tolerance stackup calculator (Python); ASME Y14.5-2018 summary; draft \texttt{gdt-annotator} SKILL.md; draft \texttt{tolerance-calculator} SKILL.md.

\textbf{Track B --- Technical Documentation and BOM} (EXEC-B, Sonnet, $\sim$25K tokens)\\
Input: BOM JSON schema (W0); title block template (W0). Output: title block generator (Python/SVG); BOM documentation; revision control pattern; draft \texttt{bom-generator} SKILL.md.

\textbf{Track C --- Vector Graphics and Programmatic Drawing} (EXEC-C, Sonnet, $\sim$38K tokens)\\
Input: Title block template (W0). Output: 20-pattern programmatic drawing cookbook (Python \texttt{svgwrite} + \texttt{ezdxf}); CNC pipeline spec; SVG-to-DXF converter; draft \texttt{svg-technical-drawer} SKILL.md; draft \texttt{blueprint-generator} SKILL.md.

\textbf{Track D --- Simulation and Testing Frameworks} (EXEC-D, Sonnet, $\sim$40K tokens)\\
Input: Materials table (from Track E via shared schema). Output: FEA tutorial (cantilever bracket); CFD conceptual tutorial; NASA-STD-7009B worksheet template; credibility level reference; draft \texttt{fea-pipeline} SKILL.md.

\textbf{Track E --- Materials Science Reference} (EXEC-E, Sonnet, $\sim$22K tokens)\\
Input: IT grade lookup table (W0). Output: 12-material properties table (JSON + Markdown); process capability map; fastener standards summary.

\textbf{Track F --- Prototype Development Pipeline} (EXEC-F, Sonnet, $\sim$23K tokens)\\
Input: All Track outputs (read-only, for pipeline coherence). Output: maker tool selection matrix; prototype pipeline tutorial (SVG $\to$ G-code); inspection protocol template; DIY build levels (L1--L4); draft \texttt{maker-blueprint} SKILL.md.

\subsection{Wave 2: Synthesis (Serial, Opus)}

\begin{center}
\rowcolors{2}{rowgray}{white}
\begin{tabularx}{\textwidth}{L{3.5cm}X}
\toprule
\rowcolor{primary}
\tableheader{Task} & \tableheader{Opus Responsibility} \\
\midrule
Career pathway synthesis & Four pathways (Mechanical Draftsman, Simulation Engineer, Maker, Technical Writer); BLS salary data; certification mapping; GSD skill mapping \\
ASME/ISO matrix & 10 high-impact differences; example drawing callouts for each \\
Try Sessions (6) & One per module; household materials feasibility review; safety check \\
Tolerance math review & Verify all stackup examples; confirm FEA reference case (analytical $\pm$2\%) \\
NASA-STD-7009B CAS & Complete Level 2 worked example for the cantilever FEA \\
Integration & Cross-module linking; through-line verification; safety review \\
\bottomrule
\end{tabularx}
\end{center}

\subsection{Wave 3: Publication and Verification (Serial, Sonnet)}

Deploy all seven SKILL.md files to \texttt{/mnt/skills/user/}. Write CK YAML entries for engineering department. Generate \texttt{index.html}. Final compilation. Release tag \texttt{v1.50.DRAFTSMAN.DRW}.

\subsection{Cache Optimization Strategy}

\begin{itemize}[noitemsep]
  \item Wave 0 must complete in $<$5 min; all Wave 1 agents receive W0 outputs in their context (cache-warm)
  \item Track C (Drawing Cookbook) is the longest track ($\sim$38K); start it first within W1
  \item Track E (Materials) is shortest; EXEC-E finishes early and deposits materials JSON for Track D to reference
  \item Wave 2 Opus receives all six Track outputs as a single assembled context (one Opus window)
  \item Wave 3 Sonnet publication is linear: SKILL deploy $\to$ CK YAML $\to$ index.html $\to$ compile
\end{itemize}

\subsection{Token Budget}

\begin{center}
\rowcolors{2}{rowgray}{white}
\begin{tabularx}{\textwidth}{L{4cm}C{2cm}C{2cm}C{2cm}X}
\toprule
\rowcolor{primary}
\tableheader{Component} & \tableheader{Model} & \tableheader{Tokens} & \tableheader{Ctx Win} & \tableheader{Notes} \\
\midrule
W0 Foundation & Haiku & 15K & 1 & Schema + scaffolding \\
Track A GD\&T & Sonnet & 32K & 2 & 14 symbols + calculator \\
Track B Docs/BOM & Sonnet & 25K & 2 & Title block + BOM \\
Track C Vector/CNC & Sonnet & 38K & 3 & 20 patterns + CNC \\
Track D FEA/Sim & Sonnet & 40K & 3 & FEA + NASA-7009B \\
Track E Materials & Sonnet & 22K & 1 & 12 materials table \\
Track F Prototype & Sonnet & 23K & 2 & L1--L4 builds \\
W2 Opus Synthesis & Opus & 45K & 2 & Careers + integration \\
W3 Publication & Sonnet & 40K & 2 & 7 SKILLs + CK + HTML \\
\midrule
\textbf{Total} & & \textbf{$\sim$280K} & \textbf{$\sim$18} & Fleet profile \\
\bottomrule
\end{tabularx}
\end{center}

\subsection{Dependency Graph}

\begin{asciibox}
W0: Shared schemas + templates
      |
      +----+----+----+----+----+----+
      |    |    |    |    |    |    |
     [A]  [B]  [C]  [D]  [E] [F]   <-- W1 parallel (6 tracks)
      |    |    |    |    |    |
      +----+----+----+----+----+
                   |
                  W2: Opus Synthesis (all 6 tracks complete)
                   |
                  W3: Publish (serial)
                   |
              v1.50.DRAFTSMAN.DRW
\end{asciibox}

\section{Test Plan}

\subsection{Test Categories}

\begin{center}
\rowcolors{2}{rowgray}{white}
\begin{tabularx}{\textwidth}{L{3.5cm}C{1.5cm}L{2cm}X}
\toprule
\rowcolor{primary}
\tableheader{Category} & \tableheader{Count} & \tableheader{Priority} & \tableheader{On Failure} \\
\midrule
Safety-Critical (math/science) & 8 & P0 & \blocktag \\
Drawing Standards Compliance & 12 & P0 & \blocktag \\
SKILL.md Schema Validation & 7 & P0 & \blocktag \\
FEA Convergence & 4 & P0 & \gatetag \\
Programmatic Drawing Output & 10 & P1 & \gatetag \\
BOM / Documentation & 6 & P1 & \annotag \\
Career Data Accuracy & 6 & P1 & \annotag \\
Try Session Feasibility & 6 & P1 & \annotag \\
Integration / Cross-Module & 8 & P2 & \annotag \\
\midrule
\textbf{Total} & \textbf{67} & & \\
\bottomrule
\end{tabularx}
\end{center}

\subsection{Safety-Critical Tests (BLOCK)}

\begin{center}
\rowcolors{2}{rowgray}{white}
\begin{tabularx}{\textwidth}{C{1.2cm}L{4cm}X}
\toprule
\rowcolor{primary}
\tableheader{ID} & \tableheader{Verifies} & \tableheader{Expected Outcome} \\
\midrule
T-001 & Cantilever FEA tip deflection & Within 2\% of $\delta = 0.800$ mm analytical \\
T-002 & Cantilever FEA peak stress & Within 2\% of $\sigma_{max} = 60.0$ MPa analytical \\
T-003 & Tolerance stackup algebraic & $T_{assy} = \sum |c_i t_i|$; verified against 3 reference cases \\
T-004 & Tolerance stackup RSS & $T_{assy} = \sqrt{\sum (c_i t_i)^2}$; verified against 3 reference cases \\
T-005 & GD\&T symbol count & All 14 characteristics present (excluding concentricity and symmetry per 2018) \\
T-006 & Process vs.\ IT grade pairing & No drawing specifies IT grade tighter than process capability \\
T-007 & Material properties sourced & Every material property value in the table has a cited source \\
T-008 & NASA-STD-7009B credibility stated & FEA tutorial explicitly declares credibility level (target: Level 2) \\
\bottomrule
\end{tabularx}
\end{center}

\subsection{Verification Matrix}

\begin{center}
\rowcolors{2}{rowgray}{white}
\begin{tabularx}{\textwidth}{L{4.5cm}L{5cm}C{2cm}}
\toprule
\rowcolor{primary}
\tableheader{Success Criterion} & \tableheader{Test IDs} & \tableheader{Status} \\
\midrule
SC1: 14 GD\&T symbols documented & T-005, T-101, T-102 & Pre-flight \\
SC2: Stackup calculator correct & T-003, T-004, T-103 & Pre-flight \\
SC3: Programmatic drawing pipeline & T-201--T-205 & Pre-flight \\
SC4: DXF importable by FreeCAD & T-206, T-207 & Pre-flight \\
SC5: FEA convergence $\pm$2\% & T-001, T-002, T-301 & Pre-flight \\
SC6: NASA-7009B worksheet complete & T-008, T-302 & Pre-flight \\
SC7: Materials table sourced & T-007, T-401 & Pre-flight \\
SC8: Prototype pipeline G-code & T-501--T-503 & Pre-flight \\
SC9: 7 SKILL.md deployed & T-601--T-607 & Pre-flight \\
SC10: Try Sessions feasible & T-701--T-706 & Pre-flight \\
SC11: ASME/ISO matrix (10 items) & T-801, T-802 & Pre-flight \\
SC12: 4 career pathways (BLS data) & T-901--T-904 & Pre-flight \\
\bottomrule
\end{tabularx}
\end{center}

\section{Mission Crew Manifest}

\begin{center}
\rowcolors{2}{rowgray}{white}
\begin{tabularx}{\textwidth}{L{2cm}L{2.8cm}C{1.5cm}X}
\toprule
\rowcolor{primary}
\tableheader{Role} & \tableheader{Agent Name} & \tableheader{Model} & \tableheader{Primary Responsibility} \\
\midrule
FLIGHT & RAVEN & Opus & Mission authority; GD\&T correctness; NASA-7009B final review \\
PLAN & SALMON & Sonnet & Wave decomposition; Track scheduling; dependency management \\
CRAFT & OWL & Opus & Career pathways; Try Sessions; integration narrative \\
EXEC-A & CEDAR & Sonnet & Drawing standards and GD\&T module \\
EXEC-B & HEMLOCK & Sonnet & Technical documentation and BOM module \\
EXEC-C & ALDER & Sonnet & Vector graphics and programmatic drawing module \\
EXEC-D & FIR & Sonnet & Simulation and testing frameworks module \\
EXEC-E & SPRUCE & Sonnet & Materials science reference module \\
EXEC-F & WILLOW & Sonnet & Prototype development pipeline module \\
TEST & HAWK & Sonnet & Test execution; convergence verification; stackup checks \\
INTEG & HERON & Sonnet & Cross-module linking; through-line verification \\
SAFETY & FOXGLOVE & Opus & Safety warden (tolerance, materials, DIY activities) \\
SKILL & OSPREY & Sonnet & 7 SKILL.md generation and schema validation \\
CK & MARTEN & Sonnet & CK YAML entries; College of Knowledge deposit \\
BUDGET & WREN & Haiku & Token tracking; context window estimates \\
SCAFFOLD & KINGFISHER & Haiku & W0 schemas; templates; config files \\
CAPCOM & [HUMAN] & --- & CAPCOM gate at W2 end; v1.50.DRAFTSMAN.DRW release approval \\
\bottomrule
\end{tabularx}
\end{center}

\section{Execution Summary}

\begin{center}
\rowcolors{2}{rowgray}{white}
\begin{tabularx}{\textwidth}{L{4cm}X}
\toprule
\rowcolor{primary}
\tableheader{Metric} & \tableheader{Value} \\
\midrule
Mission ID & v1.50.DRAFTSMAN.DRW \\
Activation profile & Fleet (17 roles) \\
Research modules & 6 (parallel, Wave 1) \\
SKILL.md deployments & 7 (\texttt{svg-technical-drawer}, \texttt{blueprint-generator}, \texttt{gdt-annotator}, \texttt{bom-generator}, \texttt{tolerance-calculator}, \texttt{fea-pipeline}, \texttt{maker-blueprint}) \\
Drawing standards covered & ASME Y14.5-2018, ISO 128-2:2020, ISO 128-3, ISO 1101:2017, ISO 2768, ISO 286-1, ANSI Y14.1 \\
Simulation frameworks & FreeCAD FEM, CalculiX, OpenFOAM, NASA-STD-7009B, ASME V\&V 40 \\
Open-source tools documented & FreeCAD, LibreCAD, OpenSCAD, CadQuery, svgwrite, ezdxf, PyCAM, dxf2gcode \\
Materials documented & 12 engineering materials $\times$ 10 properties (sourced) \\
Prototype build levels & L1 \$15 (hand tools) -- L4 \$600 (CNC router) \\
Career pathways & 4 (Mech. Draftsman, Simulation Engineer, Maker/Fabricator, Technical Writer) \\
Total tests & 67 (8 safety-critical / BLOCK) \\
CAPCOM gate & After W2 Opus synthesis, before W3 publication \\
Est. tokens & $\approx$280K \\
Est. context windows & $\approx$18 \\
Est. Opus allocation & 30\% (math review, career synthesis, safety warden) \\
Est. Sonnet allocation & 60\% (six module tracks, publication) \\
Est. Haiku allocation & 10\% (scaffolding, schemas, token tracking) \\
Release tag & \texttt{v1.50.DRAFTSMAN.DRW} on gsd-skill-creator dev branch \\
\bottomrule
\end{tabularx}
\end{center}

\begin{quotebox}
\textbf{\sffamily\large The through-line:}\\[0.5em]
{\itshape ``A drawing is not a picture of a thing. A drawing is a complete and unambiguous specification of a thing --- every geometric relationship, every allowable variation, every material property that determines whether it will work. The draftsman's art is compression: making the infinite detail of physical reality fit on a sheet of paper, in a format that any trained person on Earth can read without asking a single clarifying question.\\[0.5em]
The GSD ecosystem exists to give people their lives back. The draftsman's precision is how we make the first thing right, so we do not have to build it twice.''}\\[0.5em]
\textcolor{subtlegray}{\small --- So You Want to Be a Draftsman, v1.50.DRAFTSMAN.DRW}
\end{quotebox}

\end{document}
